\documentclass{tufte-book}

\hypersetup{colorlinks}% uncomment this line if you prefer colored hyperlinks (e.g., for onscreen viewing)
% to compile run
% pdflatex nutr630-notes
% bibtex nutr630-notes
% makeindex nutr630-notes 
% makeindex nutr630-notes.nlo -s nomencl.ist -o nutr630-notes.nls
% pdflatex nutr630-notes
% pdflatex nutr630-notes

% Book metadata
\title{Principles of Nutrition Science}
\author[Dave Bridges and Olivia S. Anderson]{Dave Bridges and Olivia Anderson}
\publisher{Department of Nutritional Sciences, University of Michigan School of Public Health}

%%
% If they're installed, use Bergamo and Chantilly from www.fontsite.com.
% They're clones of Bembo and Gill Sans, respectively.
%\IfFileExists{bergamo.sty}{\usepackage[osf]{bergamo}}{}% Bembo
%\IfFileExists{chantill.sty}{\usepackage{chantill}}{}% Gill Sans

%\usepackage{microtype}

%%
% Just some sample text
\usepackage{lipsum}
\usepackage{blindtext}
\usepackage{amsmath}

%%
% For nicely typeset tabular material
\usepackage{booktabs}

%%
% For graphics / images
\usepackage{graphicx}
\setkeys{Gin}{width=\linewidth,totalheight=\textheight,keepaspectratio}
\graphicspath{{tex-processed/figures/}}

% The fancyvrb package lets us customize the formatting of verbatim
% environments.  We use a slightly smaller font.
\usepackage{fancyvrb}
\fvset{fontsize=\normalsize}

%%
% Prints argument within hanging parentheses (i.e., parentheses that take
% up no horizontal space).  Useful in tabular environments.
\newcommand{\hangp}[1]{\makebox[0pt][r]{(}#1\makebox[0pt][l]{)}}

%%
% Prints an asterisk that takes up no horizontal space.
% Useful in tabular environments.
\newcommand{\hangstar}{\makebox[0pt][l]{*}}

%%
% Prints a trailing space in a smart way.
\usepackage{xspace}

%%
% Some shortcuts for Tufte's book titles.  The lowercase commands will
% produce the initials of the book title in italics.  The all-caps commands
% will print out the full title of the book in italics.
\newcommand{\vdqi}{\textit{VDQI}\xspace}
\newcommand{\ei}{\textit{EI}\xspace}
\newcommand{\ve}{\textit{VE}\xspace}
\newcommand{\be}{\textit{BE}\xspace}
\newcommand{\VDQI}{\textit{The Visual Display of Quantitative Information}\xspace}
\newcommand{\EI}{\textit{Envisioning Information}\xspace}
\newcommand{\VE}{\textit{Visual Explanations}\xspace}
\newcommand{\BE}{\textit{Beautiful Evidence}\xspace}

\newcommand{\TL}{Tufte-\LaTeX\xspace}

% Prints the month name (e.g., January) and the year (e.g., 2008)
\newcommand{\monthyear}{%
  \ifcase\month\or January\or February\or March\or April\or May\or June\or
  July\or August\or September\or October\or November\or
  December\fi\space\number\year
}


% Prints an epigraph and speaker in sans serif, all-caps type.
\newcommand{\openepigraph}[2]{%
  %\sffamily\fontsize{14}{16}\selectfont
  \begin{fullwidth}
  \sffamily\large
  \begin{doublespace}
  \noindent\allcaps{#1}\\% epigraph
  \noindent\allcaps{#2}% author
  \end{doublespace}
  \end{fullwidth}
}

% \index{Gut Microbiome}\index{Lipids!De Novo Lipogenesis}\index{Lipids!Lipid Synthesis}\index{Metabolites!Acetate}\index{Lipids!Fatty Acids!Short Chain Fatty Acids\index{Non-Alcoholic Fatty Liver Disease}erts a blank page
\newcommand{\blankpage}{\newpage\hbox{}\thispagestyle{empty}\newpage}

\usepackage{units}

% Typesets the font size, leading, and measure in the form of 10/12x26 pc.
\newcommand{\measure}[3]{#1/#2$\times$\unit[#3]{pc}}

% Macros for typesetting the documentation
\newcommand{\hlred}[1]{\textcolor{Maroon}{#1}}% prints in red
\newcommand{\hangleft}[1]{\makebox[0pt][r]{#1}}
\newcommand{\hairsp}{\hspace{1pt}}% hair space
\newcommand{\hquad}{\hskip0.5em\relax}% half quad space
\newcommand{\TODO}{\textcolor{red}{\bf TODO!}\xspace}
\newcommand{\na}{\quad--}% used in tables for N/A cells
\providecommand{\XeLaTeX}{X\lower.5ex\hbox{\kern-0.15em\reflectbox{E}}\kern-0.1em\LaTeX}
\newcommand{\tXeLaTeX}{\XeLaTeX\index{XeLaTeX@\protect\XeLaTeX}}
% \index{\texttt{\textbackslash xyz}@\hangleft{\texttt{\textbackslash}}\texttt{xyz}}
\newcommand{\tuftebs}{\symbol{'134}}% a backslash in tt type in OT1/T1
\newcommand{\doccmdnoindex}[2][]{\texttt{\tuftebs#2}}% command name -- adds backslash automatically (and doesn't add cmd to the index)
\newcommand{\doccmddef}[2][]{%
  \hlred{\texttt{\tuftebs#2}}\label{cmd:#2}%
  \ifthenelse{\isempty{#1}}%
    {% add the command to the index
      \index{#2 command@\protect\hangleft{\texttt{\tuftebs}}\texttt{#2}}% command name
    }%
    {% add the command and package to the index
      \index{#2 command@\protect\hangleft{\texttt{\tuftebs}}\texttt{#2} (\texttt{#1} package)}% command name
      \index{#1 package@\texttt{#1} package}\index{packages!#1@\texttt{#1}}% package name
    }%
}% command name -- adds backslash automatically
\newcommand{\doccmd}[2][]{%
  \texttt{\tuftebs#2}%
  \ifthenelse{\isempty{#1}}%
    {% add the command to the index
      \index{#2 command@\protect\hangleft{\texttt{\tuftebs}}\texttt{#2}}% command name
    }%
    {% add the command and package to the index
      \index{#2 command@\protect\hangleft{\texttt{\tuftebs}}\texttt{#2} (\texttt{#1} package)}% command name
      \index{#1 package@\texttt{#1} package}\index{packages!#1@\texttt{#1}}% package name
    }%
}% command name -- adds backslash automatically
\newcommand{\docopt}[1]{\ensuremath{\langle}\textrm{\textit{#1}}\ensuremath{\rangle}}% optional command argument
\newcommand{\docarg}[1]{\textrm{\textit{#1}}}% (required) command argument
\newenvironment{docspec}{\begin{quotation}\ttfamily\parskip0pt\parindent0pt\ignorespaces}{\end{quotation}}% command specification environment
\newcommand{\docenv}[1]{\texttt{#1}\index{#1 environment@\texttt{#1} environment}\index{environments!#1@\texttt{#1}}}% environment name
\newcommand{\docenvdef}[1]{\hlred{\texttt{#1}}\label{env:#1}\index{#1 environment@\texttt{#1} environment}\index{environments!#1@\texttt{#1}}}% environment name
\newcommand{\docpkg}[1]{\texttt{#1}\index{#1 package@\texttt{#1} package}\index{packages!#1@\texttt{#1}}}% package name
\newcommand{\doccls}[1]{\texttt{#1}}% document class name
\newcommand{\docclsopt}[1]{\texttt{#1}\index{#1 class option@\texttt{#1} class option}\index{class options!#1@\texttt{#1}}}% document class option name
\newcommand{\docclsoptdef}[1]{\hlred{\texttt{#1}}\label{clsopt:#1}\index{#1 class option@\texttt{#1} class option}\index{class options!#1@\texttt{#1}}}% document class option name defined
\newcommand{\docmsg}[2]{\bigskip\begin{fullwidth}\noindent\ttfamily#1\end{fullwidth}\medskip\par\noindent#2}
\newcommand{\docfilehook}[2]{\texttt{#1}\index{file hooks!#2}\index{#1@\texttt{#1}}}
\newcommand{\doccounter}[1]{\texttt{#1}\index{#1 counter@\texttt{#1} counter}}

% Generates the index
\usepackage{makeidx}
\makeindex

\usepackage{subfiles}
\providecommand{\alttext}[1]{}
\usepackage[intoc,norefeq]{nomencl}
\makenomenclature
% Each chapter has its own equation counter, so \nomeqref{\theequation}
% writes a different number per chapter.  Identical abbreviations from
% different chapters therefore produce different display strings, preventing
% makeindex from merging them into a single entry.  Fix: patch the internal
% writer to always use \nomeqref{0} so all copies of the same abbreviation
% are byte-for-byte identical and makeindex collapses them into one entry.
\makeatletter
\def\@@@nomenclature[#1]#2#3{%
  \def\@tempa{#2}\def\@tempb{#3}%
  \protected@write\@nomenclaturefile{}%
  {\string\nomenclatureentry{#1\nom@verb\@tempa @[{\nom@verb\@tempa}]%
      \begingroup\nom@verb\@tempb\protect\nomeqref{0}%
      |nompageref}{\thepage}}%
  \endgroup
  \@esphack}
\makeatother

\begin{document}

% Front matter
\frontmatter

% r.1 blank page
% \blankpage

% % v.2 epigraphs
% \newpage\thispagestyle{empty}
% \openepigraph{%
% The public is more familiar with bad design than good design.
% It is, in effect, conditioned to prefer bad design, 
% because that is what it lives with. 
% The new becomes threatening, the old reassuring.
% }{Paul Rand%, {\itshape Design, Form, and Chaos}
% }
% \vfill
% \openepigraph{%
% A designer knows that he has achieved perfection 
% not when there is nothing left to add, 
% but when there is nothing left to take away.
% }{Antoine de Saint-Exup\'{e}ry}
% \vfill
% \openepigraph{%
% \ldots the designer of a new system must not only be the implementor and the first 
% large-scale user; the designer should also write the first user manual\ldots 
% If I had not participated fully in all these activities, 
% literally hundreds of improvements would never have been made, 
% because I would never have thought of them or perceived 
% why they were important.
% }{Donald E. Knuth}


% r.3 full title page
\maketitle


% v.4 copyright page
\newpage
\begin{fullwidth}
~\vfill
\thispagestyle{empty}
\setlength{\parindent}{0pt}
\setlength{\parskip}{\baselineskip}
Copyright \copyright\ \the\year\ \thanklessauthor

\par\smallcaps{Published by \thanklesspublisher}

%\par\smallcaps{tufte-latex.github.io/tufte-latex/}

\par Principles of Nutritional Science by Dave Bridges and Olivia Anderson is licensed under CC BY-SA 4.0. To view a copy of this license, visit \url{https://creativecommons.org/licenses/by-sa/4.0}.

\par\textit{First printing, \monthyear}
\end{fullwidth}

% r.5 contents
\tableofcontents

\listoffigures

\listoftables

% r.7 dedication
% \cleardoublepage
% ~\vfill
% \begin{doublespace}
% \noindent\fontsize{18}{22}\selectfont\itshape
% \nohyphenation
% Dedicated to those who appreciate \LaTeX{} 
% and the work of \mbox{Edward R.~Tufte} 
% and \mbox{Donald E.~Knuth}.
% \end{doublespace}
% \vfill
% \vfill


% r.9 introduction
\cleardoublepage
\chapter*{Introduction}\label{ch:introduction}

These chapters are the lecture notes from NUTR630: Principles of Nutritional Science taught by Drs. Dave Bridges and Olivia Anderson.  Each chapter represents in general, one lecture in this course.


%%
% Start the main matter (normal chapters)
\mainmatter

\part{Regulation of Metabolism and Overview of Digestion}
\chapter{Metabolic Control Systems}\label{ch:metabolic-control-systems}
\subfile{tex-processed/metabolic-control-systems-book}

\chapter{Endocrine Control of Metabolism}\label{ch:endocrine-control}
\subfile{tex-processed/endocrine-handout-book}

\chapter{Energy Balance and Appetite}\label{ch:energy-balance}
\subfile{tex-processed/energy-balance-book}

\chapter{Overview of the Digestive System}\label{ch:intro-digestion}
\subfile{tex-processed/digestive-tract-introduction-book}

\part{Carbohydrates}

\chapter{Carbohydrate Structure}\label{ch:carb-structure}
\subfile{tex-processed/carb-structure-book}

\chapter{Carbohydrate Digestion and Transport}\label{ch:carb-digestion}
\subfile{tex-processed/carb-digestion-book}

\chapter{Fiber}\label{ch:fiber}
\subfile{tex-processed/fiber-book}

\chapter{The Gut Microbiome}\label{ch:microbiome}
\subfile{tex-processed/microbiome-book}

\chapter{Glucose Transport and Glycolysis}\label{ch:glycolysis}
\subfile{tex-processed/glycolysis-book}

\chapter{Glucose Oxidation and the TCA cycle}\label{ch:tca-cycle}
\subfile{tex-processed/tca-cycle-book}

\chapter{Oxidative Stress, the Pentose Phosphate Pathway and Alcohol Metabolism}\label{ch:ppp-ethanol}
\subfile{tex-processed/pentose-phosphate-pathway-book}

\chapter{Glycogen and Metabolism During Exercise}\label{ch:glycogen}
\subfile{tex-processed/glycogen-metabolism-book}

\chapter{Gluconeogenesis and Chronic Fasting}\label{ch:gng}
\subfile{tex-processed/gluconeogenesis-book}

\part{Lipids}

\chapter{Introduction to Lipids}\label{ch:lipid-intro}
\subfile{tex-processed/lipids-introduction-book}

\chapter{Digestion and Absorption of Lipids}\label{ch:lipid-digestion}
\subfile{tex-processed/lipid-digestion-book}

\chapter{Lipid and Cholesterol Synthesis}\label{ch:lipid-synthesis}
\subfile{tex-processed/lipid-synthesis-book}

\chapter{Lipolysis and Lipid Oxidation}\label{ch:lipid-oxidation}
\subfile{tex-processed/lipid-catabolism-book}

\chapter{Transport of Lipids and Lipoproteins}\label{ch:lipid-transport}
\subfile{tex-processed/lipid-transport-book}

\part{Proteins and Nitrogenous Compounds}

\chapter{Introduction to Proteins and Amino Acids}\label{ch:intro-proteins}
\subfile{tex-processed/proteins-amino-acids-overview-book}

\chapter{Protein Digestion and Absorption}\label{ch:protein-digestion}
\subfile{tex-processed/protein-digestion-book}

\chapter{Protein and Amino Acid Synthesis}\label{ch:protein-synthesis}
\subfile{tex-processed/proteins-amino-acids-synthesis-book}

\chapter{Protein Breakdown and Amino Acid Oxidation}\label{ch:protein-oxidation}
\subfile{tex-processed/protein-oxidation-book}

\chapter{Non-Protein Nitrogen Containing Compounds}
\label{ch:nitrogen-containing-compounds}
\subfile{tex-processed/nitrogen-compounds-book}


\backmatter
\part{Abbreviations, References and Index}
\bibliography{tex/library}
\bibliographystyle{plainnat}

\renewcommand{\nomname}{List of Abbreviations}
% Nomenclature printing for tufte-book has two problems:
%
% 1. "Lonely \item": tufte-book's \chapter* (called inside \thenomenclature)
%    creates TeX groups; \item defined with \def would be restored to the
%    original list \item inside those groups.  Fix: use \gdef.
%
% 2. "Extra \endgroup": \nompageref always ends with \endgroup (closing the
%    \begingroup that .nls entries open before the term text).  When an
%    abbreviation appears on multiple pages, makeindex emits multiple
%    \nompageref calls per entry but only one \begingroup -- every call after
%    the first has an unmatched \endgroup that unwinds outer groups, eventually
%    causing "\begin{document} ended by \end{thenomenclature}".
%    Fix: redefine \nompageref inside the environment to omit \endgroup, and
%    instead close each entry's \begingroup lazily at the start of the next
%    \item (or at \endthenomenclature for the last entry).
\makeatletter
\newlength{\@nom@labelwidth}
\setlength{\@nom@labelwidth}{3.5cm}

% Flag: are we inside the nomenclature list?
\newif\if@innomenclature
\@innomenclaturefalse

% Flag: does the preceding .nls entry have an unclosed \begingroup?
\newif\if@nom@entryopen
\@nom@entryopenfalse

% Nomenclature item handler -- [#1] matches \item [{KEY}] in .nls.
% Closes the previous entry's \begingroup before starting a new entry.
\def\@nom@nomitem[#1]{%
  \if@nom@entryopen
    \endgroup
    \global\@nom@entryopenfalse
  \fi
  \global\@nom@entryopentrue   % the .nls \begingroup for this entry follows
  \par\noindent
  \hangindent\@nom@labelwidth
  \makebox[\@nom@labelwidth][l]{\textbf{#1}}\ignorespaces
}

% Save the original \item and replace it globally with a flag-checking wrapper.
% Must use \gdef so the definition survives tufte-book's group boundaries.
\let\@nom@origitem\item
\gdef\item{%
  \if@innomenclature
    \expandafter\@nom@nomitem
  \else
    \expandafter\@nom@origitem
  \fi
}

% Save original \nompageref so we can restore it afterwards.
\let\@nom@origpageref\nompageref

% thenomenclature: add chapter heading, THEN set flag and patch \nompageref.
% The flag must be set AFTER \chapter* returns: titlesec's chapter formatting
% calls \item internally, and if our override is active during that call it
% injects \noindent/\makebox into the format argument, causing a titlesec
% "Entered in horizontal mode" error.
\renewcommand{\thenomenclature}{%
  \chapter*{\nomname}%        % chapter heading while flag is still false
  \global\@innomenclaturetrue % set AFTER \chapter* so its \item calls are unaffected
  \gdef\nompageref##1{%
    \if@printpageref\pagedeclaration{##1}\fi
    \nomentryend
    % \endgroup intentionally omitted -- see \@nom@nomitem / \endthenomenclature
  }%
}

\renewcommand{\endthenomenclature}{%
  % Close the last entry's \begingroup (no following \item to do it).
  \if@nom@entryopen
    \endgroup
    \global\@nom@entryopenfalse
  \fi
  \global\@innomenclaturefalse
  \global\let\nompageref\@nom@origpageref
  \par
}
\makeatother
\printnomenclature

\printindex

\end{document}
